% ============================================================
% 02_introduccion.tex
% ============================================================

\section{Introducción}

La gestión eficiente de los recursos energéticos en edificios educativos representa
uno de los desafíos centrales del paradigma de Edificios Inteligentes (\textit{Smart
Buildings}). Según datos de la Agencia Internacional de Energía, el sector de
edificaciones es responsable de aproximadamente el 30\,\% del consumo energético
mundial, y la iluminación artificial constituye una fracción significativa de dicho
consumo en ambientes de uso discontinuo como aulas universitarias \parencite{IEA2023}.
En estos espacios, la iluminación permanece frecuentemente encendida en ausencia
de ocupantes o cuando el nivel de luz natural resulta suficiente, lo que genera un
desperdicio energético evitable.

Los sistemas de control inteligente de iluminación han demostrado reducciones de
consumo del orden del 30\,\% al 60\,\% respecto a instalaciones convencionales
\parencite{Kralikova2016}. Sin embargo, las soluciones comerciales disponibles
presentan elevados costos de adquisición e instalación, escasa flexibilidad para
entornos de investigación y, en la mayoría de los casos, dependencia de plataformas
propietarias que dificultan su estudio y modificación.

El surgimiento de microcontroladores de bajo costo con conectividad Wi-Fi integrada
—en particular la familia ESP32 de Espressif Systems— y de plataformas IoT de código
abierto como Node-RED, han abierto la posibilidad de construir sistemas de monitoreo
y control ambiental de bajo costo, alta adaptabilidad y arquitectura abierta
\parencite{NodeRED2024, Espressif2024}. El protocolo MQTT (\textit{Message Queuing
Telemetry Transport}), diseñado originalmente por IBM para redes de baja latencia y
ancho de banda reducido, se ha consolidado como el estándar de facto para la
comunicación entre dispositivos IoT \parencite{Banks2019}.

En este contexto, el presente trabajo describe el diseño, desarrollo e implementación
de \textbf{NodeAlert IoT}, un sistema prototipo que integra:

\begin{enumerate}[label=(\roman*)]
  \item Un nodo sensor embebido basado en ESP32-C3 que mide luminosidad,
        temperatura, humedad y detecta presencia de personas mediante radar
        de microondas, aplicando una lógica local de control de iluminación
        independiente de la conectividad de red;
  \item Un servidor central desplegado sobre Raspberry Pi mediante contenedores
        Docker, que implementa un broker MQTT Mosquitto, un motor de flujos
        Node-RED con almacenamiento SQLite y una API REST;
  \item Una interfaz web desarrollada en React que permite visualizar en tiempo
        real el estado de los nodos y enviar comandos remotos.
\end{enumerate}

El trabajo forma parte del Trabajo Final de la materia Ingeniería en Computación II
de la Universidad Nacional de Rafaela (UNRaf), cuyos requisitos mínimos establecen: programación de microcontrolador con
lógica de muestreo y optimización de recursos, modos de ahorro de energía,
comunicación bidireccional por MQTT, configuración propia de broker de datos,
suscriptores y publicadores, e interfaz de usuario con visualización y controles
\parencite{CatedraICII2026}.

El artículo se organiza de la siguiente manera: la Sección~\ref{sec:marco} presenta
los fundamentos teóricos y antecedentes; la Sección~\ref{sec:materiales} detalla los
materiales y métodos; las Secciones~\ref{sec:arquitectura}
y~\ref{sec:desarrollo} describen el diseño e implementación del sistema; la
Sección~\ref{sec:resultados} presenta los resultados de validación; las
Secciones~\ref{sec:discusion} y~\ref{sec:conclusiones} ofrecen discusión y
conclusiones.

